% GENERATED FILE. Edit canonical lessons, then run npm run generate:books.
% canonical-source-hash: fnv1a64:77c7420e7c051853
% canonical-lessons: FA-C360-aqide, FA-C360-didgah, FA-C360-negaresh, FA-C360-motaqed-budan, FA-C360-nazariye

\chapter{Opinions and beliefs}
\label{ch:fa-opinions-and-beliefs}
\hlchaptermodality{360}

\begin{chapteropening}
\textbf{By the end of this chapter:} \emph{I can give my opinion and my point of view, say what I believe and name a theory.}

Complete the last lesson of chapter 360: I can give my opinion and my point of view, say what I believe and name a theory.
\end{chapteropening}

\section[\texorpdfstring{\fa{عقیده} ('aqide)}{عقیده ('aqide)}]{\fa{عقیده} ('aqide) — an opinion, a belief}

\label{lesson:FA-C360-aqide}



\begin{quote}
Before the new one: say the Persian for an invoice, then the Persian for a cash machine.
\end{quote}

\subsection*{You'll want to know: \fa{عقیده}}
\textbf{\fa{عقیده}} (\emph{'aqide}) — \textquotedblleft{}an opinion, a belief\textquotedblright{}.

\begin{quote}
\fa{عقیده} \fa{من} \fa{این} \fa{است} \fa{که} \fa{او} \fa{راست} \fa{میگوید}.
\end{quote}

\begin{quote}
\emph{'aqide-ye man in ast ke u râst miguyad.}
\end{quote}

\begin{quote}
My opinion is that he is telling the truth.
\end{quote}

Plural \emph{'aqâyed}, an Arabic broken plural. \emph{Be 'aqide-ye man} — in my opinion, a little more formal than \emph{be nazar-e man}.

\begin{cousinweb}
\textbf{\fa{عقیده}} was borrowed from Arabic \emph{'aqīda} \textquotedblleft{}a creed, a firm belief\textquotedblright{} (EtymDB). A belief, in this picture, is something you have tied fast.
\end{cousinweb}

\begin{grammarlens}[title={What you've built}]
One more word for telling, explaining and describing at B1.
\end{grammarlens}

\subsection*{Your turn}
\begin{itemize}
\raggedright
  \item \emph{Say it:} \emph{'aqide}
  \item \emph{Say it:} \emph{'aqide}, once more
  \item \emph{Say it:} a sentence of your own with \emph{'aqide}
  \item \emph{Recall:} say the Persian for an invoice, then the Persian for a cash machine, then say \emph{'aqide} again
\end{itemize}

\begin{tcolorbox}[breakable,colback=teal!4,colframe=teal!35!black,title={Before you move on}]
What does \fa{عقیده} mean? (\textquotedblleft{}An opinion, a belief\textquotedblright{}.) Say it once more.
\end{tcolorbox}
\section[\texorpdfstring{\fa{دیدگاه} (didgâh)}{دیدگاه (didgâh)}]{\fa{دیدگاه} (didgâh) — a point of view, a viewpoint}

\label{lesson:FA-C360-didgah}



\begin{quote}
Before the new one: say the Persian for a cash machine, then the Persian for an opinion.
\end{quote}

\subsection*{You'll want to know: \fa{دیدگاه}}
\textbf{\fa{دیدگاه}} (\emph{didgâh}) — \textquotedblleft{}a point of view, a viewpoint\textquotedblright{}.

\begin{quote}
\fa{از} \fa{دیدگاه} \fa{پدرم،} \fa{این} \fa{کار} \fa{خطرناک} \fa{است}.
\end{quote}

\begin{quote}
\emph{az didgâh-e pedaram, in kâr khatarnâk ast.}
\end{quote}

\begin{quote}
From my father's point of view, this is dangerous.
\end{quote}

\emph{Az didgâh-e} … — from the point of view of … . Plural \emph{didgâhhâ}.

\begin{cousinweb}
\textbf{\fa{دیدگاه}} is a native Persian compound built from \emph{did} \textquotedblleft{}sight\textquotedblright{}, the past stem of \emph{didan} \textquotedblleft{}to see\textquotedblright{}, and \emph{-gâh} \textquotedblleft{}place\textquotedblright{}. A point of view is the place you look from.
\end{cousinweb}

\begin{grammarlens}[title={What you've built}]
One more word for telling, explaining and describing at B1.
\end{grammarlens}

\subsection*{Your turn}
\begin{itemize}
\raggedright
  \item \emph{Say it:} \emph{didgâh}
  \item \emph{Say it:} \emph{didgâh}, once more
  \item \emph{Say it:} a sentence of your own with \emph{didgâh}
  \item \emph{Recall:} say the Persian for a cash machine, then the Persian for an opinion, then say \emph{didgâh} again
\end{itemize}

\begin{tcolorbox}[breakable,colback=teal!4,colframe=teal!35!black,title={Before you move on}]
What does \fa{دیدگاه} mean? (\textquotedblleft{}A point of view, a viewpoint\textquotedblright{}.) Say it once more.
\end{tcolorbox}
\section[\texorpdfstring{\fa{نگرش} (negaresh)}{نگرش (negaresh)}]{\fa{نگرش} (negaresh) — an attitude, an outlook}

\label{lesson:FA-C360-negaresh}



\begin{quote}
Before the new one: say the Persian for an opinion, then the Persian for a point of view.
\end{quote}

\subsection*{You'll want to know: \fa{نگرش}}
\textbf{\fa{نگرش}} (\emph{negaresh}) — \textquotedblleft{}an attitude, an outlook\textquotedblright{}.

\begin{quote}
\fa{نگرش} \fa{او} \fa{به} \fa{کار} \fa{عوض} \fa{شده} \fa{است}.
\end{quote}

\begin{quote}
\emph{negaresh-e u be kâr 'avaz shode ast.}
\end{quote}

\begin{quote}
His attitude to work has changed.
\end{quote}

\emph{Negaresh be} — an attitude towards. Formal and written; in speech people often say \emph{tarz-e fekr}.

\begin{cousinweb}
\textbf{\fa{نگرش}} is built on \emph{negar-}, the present stem of the native verb \emph{negaristan} \textquotedblleft{}to look\textquotedblright{} (EtymDB lists the verb), with the noun suffix \emph{-esh}: an attitude is a way of looking.
\end{cousinweb}

\begin{grammarlens}[title={What you've built}]
One more word for telling, explaining and describing at B1.
\end{grammarlens}

\subsection*{Your turn}
\begin{itemize}
\raggedright
  \item \emph{Say it:} \emph{negaresh}
  \item \emph{Say it:} \emph{negaresh}, once more
  \item \emph{Say it:} a sentence of your own with \emph{negaresh}
  \item \emph{Recall:} say the Persian for an opinion, then the Persian for a point of view, then say \emph{negaresh} again
\end{itemize}

\begin{tcolorbox}[breakable,colback=teal!4,colframe=teal!35!black,title={Before you move on}]
What does \fa{نگرش} mean? (\textquotedblleft{}An attitude, an outlook\textquotedblright{}.) Say it once more.
\end{tcolorbox}
\section[\texorpdfstring{\fa{معتقد} \fa{بودن} (mo'taqed budan)}{معتقد بودن (mo'taqed budan)}]{\fa{معتقد} \fa{بودن} (mo'taqed budan) — to believe, to be convinced}

\label{lesson:FA-C360-motaqed-budan}



\begin{quote}
Before the new one: say the Persian for a point of view, then the Persian for an attitude.
\end{quote}

\subsection*{You'll want to know: \fa{معتقد} \fa{بودن}}
\textbf{\fa{معتقد} \fa{بودن}} (\emph{mo'taqed budan}) — \textquotedblleft{}to believe, to be convinced\textquotedblright{}.

\begin{quote}
\fa{من} \fa{معتقدم} \fa{که} \fa{ورزش} \fa{برای} \fa{همه} \fa{لازم} \fa{است}.
\end{quote}

\begin{quote}
\emph{man mo'taqedam ke varzesh barâye hame lâzem ast.}
\end{quote}

\begin{quote}
I believe that sport is necessary for everyone.
\end{quote}

\emph{Mo'taqed} \textquotedblleft{}convinced\textquotedblright{} + \emph{budan} \textquotedblleft{}to be\textquotedblright{}: \emph{mo'taqedam ke} … — I believe that … . A belief in something takes \emph{be}: \emph{be chizi mo'taqed budan}.

\begin{cousinweb}
\textbf{\fa{معتقد}} is usually traced to Arabic \emph{mu'taqid} \textquotedblleft{}believing, holding firmly\textquotedblright{}, from the same Arabic family as \emph{'aqīda}. \emph{Budan} \textquotedblleft{}to be\textquotedblright{} is the inherited Persian verb.
\end{cousinweb}

\begin{grammarlens}[title={What you've built}]
One more word for telling, explaining and describing at B1.
\end{grammarlens}

\subsection*{Your turn}
\begin{itemize}
\raggedright
  \item \emph{Say it:} \emph{mo'taqed budan}
  \item \emph{Say it:} \emph{mo'taqed budan}, once more
  \item \emph{Say it:} a sentence of your own with \emph{mo'taqed budan}
  \item \emph{Recall:} say the Persian for a point of view, then the Persian for an attitude, then say \emph{mo'taqed budan} again
\end{itemize}

\begin{tcolorbox}[breakable,colback=teal!4,colframe=teal!35!black,title={Before you move on}]
What does \fa{معتقد} \fa{بودن} mean? (\textquotedblleft{}To believe, to be convinced\textquotedblright{}.) Say it once more.
\end{tcolorbox}
\section[\texorpdfstring{\fa{نظریه} (nazariye)}{نظریه (nazariye)}]{\fa{نظریه} (nazariye) — a theory}

\label{lesson:FA-C360-nazariye}



\begin{quote}
Before the new one: say the Persian for an attitude, then the Persian for to believe.
\end{quote}

\subsection*{You'll want to know: \fa{نظریه}}
\textbf{\fa{نظریه}} (\emph{nazariye}) — \textquotedblleft{}a theory\textquotedblright{}.

\begin{quote}
\fa{این} \fa{نظریه} \fa{خیلی} \fa{قدیمی} \fa{است}.
\end{quote}

\begin{quote}
\emph{in nazariye kheyli qadimi ast.}
\end{quote}

\begin{quote}
This theory is very old.
\end{quote}

Built on \emph{nazar} \textquotedblleft{}view\textquotedblright{} with the suffix \emph{-iye}. Plural \emph{nazariyehâ}.

\begin{cousinweb}
\textbf{\fa{نظریه}} comes from Arabic \emph{naẓariyya} \textquotedblleft{}theory\textquotedblright{} (EtymDB), built on \emph{naẓar} \textquotedblleft{}a look, a view\textquotedblright{}, the same word that gave Persian \emph{nazar}. A theory is a way of viewing the facts.
\end{cousinweb}

\begin{grammarlens}[title={What you've built}]
One more word for telling, explaining and describing at B1.
\end{grammarlens}

\subsection*{Your turn}
\begin{itemize}
\raggedright
  \item \emph{Say it:} \emph{nazariye}
  \item \emph{Say it:} \emph{nazariye}, once more
  \item \emph{Say it:} a sentence of your own with \emph{nazariye}
  \item \emph{Recall:} say the Persian for an attitude, then the Persian for to believe, then say \emph{nazariye} again
\end{itemize}

\begin{tcolorbox}[breakable,colback=teal!4,colframe=teal!35!black,title={Before you move on}]
What does \fa{نظریه} mean? (\textquotedblleft{}A theory\textquotedblright{}.) Say it once more.
\end{tcolorbox}
